We assess how connection loss affects collision avoidance in a fly-connectome spiking neural network (SNN). The synthetic pursuit--escape task uses paired deletion masks and two frozen feedback laws. In the baseline receptive-field realization (RF-0), both laws complete all 48 healthy tasks, yet 50\% connection-row deletion yields 37/48 successes under proportional control and 48/48 under relay control. The prespecified secondary relay-minus-proportional net-loss contrast is $-22.9$ percentage points (95\% whole-seed CI $[-33.3,-12.5]$); two independently recalibrated same-prior realizations each yield an uncertain $-4.2$ $[-12.5,0]$. Recalibration, different test cohorts and unequal healthy clearance qualify this comparison. A frozen command-rate warning detects 2/13, 2/3 and 0/1 subsequent collisions in RF-0/A/B. An executable evidence interface checks pairing, timing and task labels, with independent native replay and demonstrated reuse on a delayed tank. Together, these results support task-level assessment for a declared sensing--control realization and fixed 1-s collision horizon.
